% ============================================================
%  DigiRadio — Manual chapter: FSC-BT1035 Bluetooth
% ============================================================

\chapter{FSC-BT1035 Bluetooth Transmitter}
\label{ch:bt1035}

The Feasycom FSC-BT1035 (Qualcomm QCC3056) is the wireless output stage of
DigiRadio: it receives PCM from the ADAU1701 over I\textsuperscript{2}S and
streams Bluetooth audio with aptX, aptX~HD, and aptX~Adaptive. This chapter
documents how the ESP32-S3 controls the module over UART (AT commands, no
hardware flow control), why I\textsuperscript{2}S slave mode is mandatory,
and how \texttt{bt1035::Bt1035Driver} implements the bring-up sequence.

\begin{drref}[Hardware context]
Board wiring (UART pins, I\textsuperscript{2}S to the module, flow control)
is in Chapter~\ref{ch:hardware}, Section~\ref{sec:hw-bt1035}. The ADAU1701
master clock and limiter settings that feed the module are in
Chapter~\ref{ch:adau1701} and Chapter~\ref{ch:sigmastudio}.
\end{drref}

\section{Role in the audio chain}
\label{sec:bt1035-role}

The BT1035 is an I\textsuperscript{2}S \textbf{slave} at 48\,kHz: LRCLK and
BCLK come from the ADAU1701; PCM arrives on \texttt{SDATA\_OUT0}
(ADAU MP6 $\rightarrow$ module pin~5). The ESP32 does not process audio
samples for Bluetooth; it only configures the module so the wired path is
accepted and encoded for transmission.

Without firmware init the module may stay in a default mode that ignores the
I\textsuperscript{2}S bus from the ADAU1701. The mandatory
\texttt{AT+AUXCFG=3} and \texttt{AT+I2SCFG=35} commands select I\textsuperscript{2}S
slave input at 48\,kHz, 24-bit --- omitting them silently breaks the entire
wireless output (Section~\ref{sec:bt1035-i2s}).

\section{Control interface}
\label{sec:bt1035-uart}

\subsection{UART parameters}

\begin{table}[htbp]
  \centering
  \small
  \begin{tabular}{@{}L{4.2cm}L{9.2cm}@{}}
    \drhead Setting & Value \\
    \midrule
    Port            & UART2 (not the console UART) \\
    Baud rate       & 115200 \\
    Data format     & 8N1 \\
    Flow control    & Disabled on the ESP32 UART (\texttt{UART\_HW\_FLOWCTRL\_DISABLE}) \\
    Reset           & GPIO active-low pulse at boot \\
    SYS\_CTL        & Held active to enable the module \\
    \bottomrule
  \end{tabular}
  \caption{FSC-BT1035 UART and control lines (\texttt{board\_pins.hpp}).}
  \label{tab:bt1035-uart}
\end{table}

\begin{drnote}[Flow control disabled, not just unused]
Earlier revisions wired \texttt{RTS}/\texttt{CTS} through
\texttt{Bt1035Pins} and configured \texttt{UART\_HW\_FLOWCTRL\_CTS\_RTS}.
That was reverted: \texttt{Bt1035Driver.cpp} now configures
\texttt{UART\_HW\_FLOWCTRL\_DISABLE} and passes
\texttt{UART\_PIN\_NO\_CHANGE} for both lines --- GPIO14/21 are not driven by
this driver. No AT command in \texttt{core::Bt1035AtCommand} configures flow
control on the module side either; this has not caused observed boot
failures, but has not been independently verified against the module's own
flow-control expectation.
\end{drnote}

\subsection{Response handling}

Every command expects a module reply containing \texttt{OK} or
\texttt{ERROR}. The driver:

\begin{itemize}
  \item builds lines in the pure core via \texttt{core::buildBt1035AtLine()};
  \item classifies replies with \texttt{core::parseBt1035AtResponse()};
  \item treats timeouts and unexpected payloads as errors (never ignored).
\end{itemize}

Host tests in \drpath{components/core/test/bt1035_at_test.cpp} lock the
init sequence (\texttt{AT+AUXCFG=3}, \texttt{AT+I2SCFG=35}) and the parser.

\section{Mandatory I\textsuperscript{2}S slave mode}
\label{sec:bt1035-i2s}

\begin{drcaution}[I\textsuperscript{2}S init is not optional]
The board routes ADAU1701 \texttt{SDATA\_OUT0} (MP6) to the module PCM input
with shared BCLK/LRCLK (see Chapter~\ref{ch:hardware}). The init sequence
must use:
\begin{itemize}
  \item \texttt{AT+AUXCFG=3} --- I\textsuperscript{2}S mode (§5.1.25)
  \item \texttt{AT+I2SCFG=35} --- I\textsuperscript{2}S slave, 48\,kHz, 24-bit (§5.1.4),
    matching the ADAU1701 serial output word length configured in the
    SigmaStudio Hardware Configuration
\end{itemize}
\texttt{AT+AUXCFG=1} (Line-In) does \textbf{not} match the schematic.
AGENTS.md treats skipping I\textsuperscript{2}S init as a production bug.
\end{drcaution}

\section{Boot sequence}
\label{sec:bt1035-boot}

At power-up \texttt{HardwareBootstrap::boot()} runs the Si4684 and ADAU1701
first, applies the saved audio profile, then initialises the BT1035 so the
I\textsuperscript{2}S path is ready before Wi-Fi starts.

\begin{figure}[htbp]
  \centering
  \begin{tikzpicture}[
      drstep/.style={draw, rounded corners, minimum width=9.5cm,
        minimum height=0.9cm, align=center, font=\small},
      >={Latex}, node distance=3mm]
    \node[drstep, fill=black!6] (sys) {SYS\_CTL high, RESET\# pulse};
    \node[drstep, fill=black!6, below=of sys] (uart)
      {Install UART2 @ 115200, flow control disabled};
    \node[drstep, fill=black!7, below=of uart] (swrst)
      {Send \texttt{AT+RESET} (best-effort), settle 500\,ms, flush RX};
    \node[drstep, fill=black!8, below=of swrst] (at)
      {Send \texttt{AT} --- expect OK};
    \node[drstep, fill=black!10, below=of at] (aux)
      {Send \texttt{AT+AUXCFG=3} --- expect OK (I\textsuperscript{2}S)};
    \node[drstep, fill=black!10, below=of aux] (i2s)
      {Send \texttt{AT+I2SCFG=35} --- expect OK (slave 48\,kHz, 24-bit)};
    \node[drstep, fill=black!6, below=of i2s] (done)
      {\texttt{Bt1035Driver::isBooted()} = true};
    \foreach \a/\b in {sys/uart, uart/swrst, swrst/at, at/aux, aux/i2s, i2s/done} {
      \draw[->] (\a) -- (\b);
    }
  \end{tikzpicture}
  \caption{FSC-BT1035 AT init on DigiRadio.}
  \label{fig:bt1035-boot}
\end{figure}

\begin{drnote}[AT+RESET is best-effort, not gated]
Unlike \texttt{AT+AUXCFG=3}/\texttt{AT+I2SCFG=35} (mandatory,
\texttt{core::bootInitSequence()}, boot fails if either is refused),
\texttt{AT+RESET} is sent separately and its result is discarded: it is
unverified whether this Feasycom firmware acknowledges the command before
rebooting, resets silently without a reply, or rejects it outright. Gating
boot on an unconfirmed ack would risk breaking an otherwise-working sequence.
The 500\,ms settle delay and RX flush after sending it mirror the delay
already used after the hardware GPIO reset pulse.
\end{drnote}

Pairing, codec selection, and volume over Bluetooth are handled by the
module's own firmware and NVS; DigiRadio firmware currently implements
only the I\textsuperscript{2}S bring-up required for the wired audio path.

\section{Software architecture}
\label{sec:bt1035-stack}

\begin{table}[htbp]
  \centering
  \small
  \begin{tabular}{@{}L{2.5cm}L{3.8cm}L{6.2cm}@{}}
    \drhead Layer & Type & Responsibility \\
    \midrule
    Bootstrap     & \texttt{HardwareBootstrap} &
      Constructs driver, calls \texttt{boot()} after ADAU1701 \\
    Driver        & \texttt{bt1035::Bt1035Driver} &
      UART, reset, AT init sequence \\
    Pure core     & \texttt{core::Bt1035AtCommand}, parser &
      Host-testable command strings and OK/ERROR classification \\
    \bottomrule
  \end{tabular}
  \caption{BT1035 control stack (Slice~6--7). Pairing and A2DP status are
    exposed on \texttt{/api/bluetooth/*} via \texttt{BluetoothService}.}
  \label{tab:bt1035-stack}
\end{table}

\section{Supported AT command subset}
\label{sec:bt1035-at}

The firmware enumerates every command it sends. Wire formats follow
\drpath{Hardware/DATASHEET/FSC-BT1035_programming_user_guide_1.1.1.pdf}
(§5 commands, §6 events). Extending the subset requires updating
\texttt{core::Bt1035AtCommand}, the manual, and a host test.

\begin{table}[htbp]
  \centering
  \small
  \begin{tabular}{@{}L{2.4cm}L{3.4cm}L{5.5cm}@{}}
    \drhead Enum & Line sent & Programming guide \\
    \midrule
    \texttt{Reset}          & \texttt{AT+RESET}        & software reset, best-effort \\
    \texttt{Ping}           & \texttt{AT}              & link check \\
    \texttt{I2sMode}        & \texttt{AT+AUXCFG=3}     & §5.1.25 Param=3 I2S \\
    \texttt{I2sSlave48k24}  & \texttt{AT+I2SCFG=35}    & §5.1.4 slave 48\,kHz 24-bit \\
    \texttt{PairDiscoverable} & \texttt{AT+PAIR=1} & §5.1.20 enter discoverable \\
    \texttt{PairHidden}       & \texttt{AT+PAIR=0} & §5.1.20 leave discoverable \\
    \texttt{A2dpStat}         & \texttt{AT+A2DPSTAT} & §5.3.1; states 1--5 \\
    \texttt{A2dpDisconnect}   & \texttt{AT+A2DPDISC} & §5.3.3 \\
    \texttt{QueryName}        & \texttt{AT+NAME}     & §5.1.16 read \texttt{+NAME=} \\
    \texttt{QueryAutoConn}    & \texttt{AT+AUTOCONN} & §5.1.11 read \texttt{+AUTOCONN=} \\
    \texttt{QueryPairedList}  & \texttt{AT+PLIST}    & §5.1.22; ends with \texttt{+PLIST=E} \\
    \bottomrule
  \end{tabular}
  \caption{Enumerated AT commands (\texttt{core::Bt1035AtCommand}). Boot
    sends Reset (best-effort) then the mandatory Ping + I2sMode +
    I2sSlave48k24 (\texttt{core::bootInitSequence()}); pairing commands are
    runtime.}
  \label{tab:bt1035-at}
\end{table}

Boot also calls \texttt{AT+NAME=<identity>,0} and \texttt{AT+AUTOCONN=3}
from \texttt{hardware\_bootstrap.cpp} (§5.1.16 suffix disabled, §5.1.11
reconnect count).

\section{Bt1035Driver API}
\label{sec:bt1035-driver}

\texttt{Bt1035Driver} owns UART and GPIO reset. All methods return
\texttt{std::expected<T, Bt1035Error>}.

\begin{table}[htbp]
  \centering
  \small
  \begin{tabular}{@{}L{4.2cm}L{9.2cm}@{}}
    \toprule
    \textbf{Method} & \textbf{Purpose} \\
    \midrule
    \texttt{boot()} & Reset, UART init, run \texttt{bootInitSequence()} \\
    \texttt{isBooted()} & \texttt{true} after I\textsuperscript{2}S init succeeded \\
    \texttt{sendCommand(cmd)} & Send one typed command, expect OK \\
    \texttt{enterPairingMode()} & \texttt{AT+PAIR=1} \\
    \texttt{leavePairingMode()} & \texttt{AT+PAIR=0} \\
    \texttt{queryA2dpState()} & \texttt{AT+A2DPSTAT}, parse \texttt{+A2DPSTAT=} \\
    \texttt{disconnectA2dp()} & \texttt{AT+A2DPDISC} \\
    \texttt{queryDeviceName()} & \texttt{AT+NAME}, parse \texttt{+NAME=} \\
    \texttt{queryAutoReconnect()} & \texttt{AT+AUTOCONN}, parse \texttt{+AUTOCONN=} \\
    \texttt{setAutoReconnect(times)} & \texttt{AT+AUTOCONN=n} (0--15) \\
    \texttt{queryPairedList()} & \texttt{AT+PLIST}, parse \texttt{+PLIST=} lines \\
    \bottomrule
  \end{tabular}
  \caption{Public driver API.}
  \label{tab:bt1035-api}
\end{table}

\subsection{Error codes}

\begin{table}[htbp]
  \centering
  \small
  \begin{tabular}{@{}L{4.2cm}L{9.2cm}@{}}
    \drhead \texttt{Bt1035Error} & Typical cause \\
    \midrule
    \texttt{ResetFailed}         & GPIO configuration failure \\
    \texttt{UartInitFailed}      & \texttt{uart\_driver\_install} / pins \\
    \texttt{NotBooted}           & \texttt{sendCommand} before \texttt{boot()} \\
    \texttt{AtTimeout}           & No OK/ERROR within 2\,s \\
    \texttt{AtError}             & Module returned ERROR \\
    \texttt{UnexpectedResponse}  & Unrecognised payload (future use) \\
    \bottomrule
  \end{tabular}
  \caption{Driver error enumeration.}
  \label{tab:bt1035-errors}
\end{table}

\section{Integration at power-up}
\label{sec:bt1035-integration}

\texttt{main/hardware\_bootstrap.cpp} constructs a static
\texttt{Bt1035Driver} and calls \texttt{boot()} after
\texttt{AudioService::loadAndApply()}. Failure returns
\texttt{HardwareBootError::Bt1035BootFailed} and \texttt{app\_main} halts
before network bring-up (fail-closed, same as Si4684/ADAU1701).

Boot order:
\begin{enumerate}
  \item Si4684 \texttt{boot(Dab)} --- tuner image in RAM.
  \item ADAU1701 \texttt{boot()} --- SigmaStudio program in RAM.
  \item \texttt{AudioService::loadAndApply()} --- user mixer/EQ profile.
  \item BT1035 \texttt{boot()} --- I\textsuperscript{2}S slave enabled for wireless output.
\end{enumerate}

\section{Typical usage (firmware developer)}
\label{sec:bt1035-usage}

After a successful \texttt{HardwareBootstrap::boot()}, the module is ready;
no further calls are required for basic listening. To re-send I\textsuperscript{2}S config
after a module reset:

\begin{verbatim}
bt1035::Bt1035Driver& bt = ...;
if (auto r = bt.sendCommand(core::Bt1035AtCommand::I2sMode); !r) {
  // handle Bt1035Error
}
if (auto r = bt.sendCommand(core::Bt1035AtCommand::I2sSlave48k24); !r) {
  // handle Bt1035Error
}
\end{verbatim}

\section{Further reading}
\label{sec:bt1035-reading}

\begin{itemize}
  \item \texttt{Hardware/DATASHEET/FSC-BT1035\_programming\_user\_guide\_1.1.1.pdf}
        --- authoritative AT command and event reference (§5--§6).
  \item \texttt{Hardware/DATASHEET/FSC-BT1035\_Datasheet\_EN.pdf} --- module
        electrical and pinout specification.
  \item Chapter~\ref{ch:hardware} --- pin map and I\textsuperscript{2}S routing.
  \item Chapter~\ref{ch:adau1701} --- DSP output that feeds the module.
  \item \texttt{components/core/test/bt1035\_at\_test.cpp} --- init sequence
        and parser host tests.
\end{itemize}
